\documentclass{amsart}
\usepackage{amsmath,amssymb,amsthm,hyperref}

\theoremstyle{plain}
\newtheorem{theorem}{Theorem}
\newtheorem{lemma}{Lemma}
\newtheorem{proposition}{Proposition}
\newtheorem{corollary}{Corollary}
\theoremstyle{definition}
\newtheorem{remark}{Remark}

\DeclareMathOperator{\rank}{rank}
\DeclareMathOperator{\spann}{span}
\DeclareMathOperator{\cl}{cl}
\newcommand{\I}{\mathcal{I}}
\newcommand{\F}{\mathbb{F}}

\begin{document}

\title{A polynomial-time test for equality of two linearly represented matroids}
\maketitle

\section*{Statement}

Fix a finite field $\F=\mathrm{GF}(q)$, $q=p^{k}$ ($p$ prime). For a matrix $A$ over
$\F$ whose columns are indexed by a finite set $E$, let $M(A)=(E,\I(A))$ be the vector
matroid: $S\subseteq E$ is independent ($S\in\I(A)$) iff the columns $\{A_e:e\in S\}$
are linearly independent over $\F$. Write $r_A(S)=\dim\spann_{\F}\{A_e:e\in S\}
=\rank_{\F}\bigl(A[\,\cdot\,,S]\bigr)$ for its rank function.

\begin{theorem}\label{thm:main}
Fix the finite field $\F$. There is an algorithm that, given two matrices $A_1,A_2$ over
$\F$ with the same column-index set $E$, decides whether $M(A_1)=M(A_2)$ in time
polynomial in the total bit-length of the input. In particular the answer to the
question is \emph{yes}.
\end{theorem}

The proof is organized as follows. Section~\ref{sec:prelim} fixes matroid preliminaries and
records the rank oracle (Lemma~\ref{lem:oracle}). Section~\ref{sec:reduction} reduces equality
to two one-sided \emph{inclusion} tests (Lemma~\ref{lem:reduction}) and characterizes inclusion
by a basis condition (Lemma~\ref{lem:basis}). Section~\ref{sec:minmax} proves the key
combinatorial min--max identity (Lemma~\ref{lem:minmax}) and deduces that the relevant quantity
is a weighted matroid-intersection optimum, hence polynomially computable
(Lemma~\ref{lem:wmi}). Section~\ref{sec:algorithm} assembles the algorithm, proves correctness,
the complexity bound, and the degenerate cases, completing the proof of
Theorem~\ref{thm:main}. Section~\ref{sec:verif} summarizes the exhaustive numerical checks of all
combinatorial claims.

Throughout, $n=|E|$ and we abbreviate $r_i:=r_{A_i}$ and $M_i:=M(A_i)$.

\section{Preliminaries and the rank oracle}\label{sec:prelim}

We use the standard facts of matroid theory (Oxley, \emph{Matroid Theory}, 2nd ed.; Schrijver,
\emph{Combinatorial Optimization}, Vol.~B).

\begin{itemize}
\item A rank function $r:2^E\to\mathbb Z_{\ge0}$ of a matroid satisfies $r(\varnothing)=0$,
monotonicity $r(S)\le r(T)$ for $S\subseteq T$, unit increase $r(S)\le r(S\cup\{e\})\le r(S)+1$,
$r(S)\le|S|$, and submodularity $r(S)+r(T)\ge r(S\cup T)+r(S\cap T)$. For a vector matroid $M(A)$,
$r_A$ is exactly the linear-algebra rank, which enjoys all of these.

\item A set $S$ is independent iff $r(S)=|S|$. A \emph{basis} is a maximal independent set;
all bases have cardinality $r(E)$. The \emph{closure} of $S$ is
$\cl_M(S)=\{e\in E: r(S\cup\{e\})=r(S)\}$.

\item Two matroids on the same ground set $E$ are equal iff they have the same family of
independent sets, equivalently iff they have the same rank function.

\item (Subadditivity used below.) For any matroid with rank $r$ and any sets $S,U$,
$r(S\cup U)\le r(S)+|U|$ (unit increase applied $|U|$ times); in particular for $S\subseteq B$,
$r(B)\le r(S)+(|B|-|S|)$.
\end{itemize}

\begin{lemma}[Rank oracle]\label{lem:oracle}
There is an algorithm that, given a matrix $A$ over $\F=\mathrm{GF}(q)$ on column set $E$ and a
subset $S\subseteq E$, outputs $r_A(S)$ in time polynomial in the bit-length of $A$.
\end{lemma}

\begin{proof}
Represent $\mathrm{GF}(p^k)$ as $\mathbb F_p[x]/(f)$ for a fixed irreducible $f$ of degree $k$
over $\mathbb F_p$; this representation is part of the fixed field, not of the input. A field
element is a polynomial of degree $<k$ with coefficients in $\{0,\dots,p-1\}$, i.e.
$O(k\log p)$ bits, which is $O(1)$ since $\F$ is fixed (and $O(\log q)$ in general). Addition,
multiplication, and inversion in $\mathrm{GF}(p^k)$ are polynomial arithmetic modulo $f$ and the
extended Euclidean algorithm, each $\mathrm{poly}(\log q)=O(1)$ bit operations. Gaussian
elimination on $A[\,\cdot\,,S]$ performs $O\!\left(\min(m,|S|)\,m\,|S|\right)$ field operations
($m=$ number of rows), polynomial in the number of entries of $A$, hence in the bit-length of
$A$. The number of nonzero pivots equals $r_A(S)$. \qedhere
\end{proof}

Thus $r_1,r_2$ are available as polynomial-time \emph{rank oracles}; every running time below is
polynomial in the input bit-length whenever it is polynomial in $n$ together with a polynomial
number of oracle calls.

\section{Reduction to a one-sided inclusion test}\label{sec:reduction}

\begin{lemma}[Reduction to two inclusions]\label{lem:reduction}
For matrices $A_1,A_2$ on $E$,
\[
   M(A_1)=M(A_2)\quad\Longleftrightarrow\quad \I(A_1)\subseteq\I(A_2)\ \text{ and }\ \I(A_2)\subseteq\I(A_1).
\]
Hence it suffices to decide a single inclusion $\I(A_1)\subseteq\I(A_2)$ in polynomial time and
to run the same procedure with $A_1,A_2$ interchanged.
\end{lemma}

\begin{proof}
Two matroids are equal iff they have the same independent sets, i.e. $\I(A_1)=\I(A_2)$, the
conjunction of the two displayed inclusions. \qedhere
\end{proof}

For a basis $B$ of $M_1$ we have $|B|=r_1(E)$ and $r_2(B)\le|B|$, with equality iff
$B\in\I(A_2)$.

\begin{lemma}[Basis characterization of inclusion]\label{lem:basis}
$\I(A_1)\subseteq\I(A_2)$ holds if and only if every basis $B$ of $M_1$ satisfies
$B\in\I(A_2)$; equivalently,
\[
   Q:=\min_{B\ \text{basis of}\ M_1} r_2(B)\;=\;r_1(E).
\]
\end{lemma}

\begin{proof}
($\Rightarrow$) Every basis $B$ of $M_1$ is $M_1$-independent, hence $M_2$-independent, so
$r_2(B)=|B|=r_1(E)$; thus $Q=r_1(E)$.

($\Leftarrow$) Suppose $\I(A_1)\not\subseteq\I(A_2)$; choose $S\in\I(A_1)\setminus\I(A_2)$. By the
augmentation axiom $S$ extends to a basis $B\supseteq S$ of $M_1$. Since $\I(A_2)$ is closed under
taking subsets and $S\notin\I(A_2)$, also $B\notin\I(A_2)$, i.e. $r_2(B)<|B|=r_1(E)$, so
$Q<r_1(E)$. Contrapositively, $Q=r_1(E)$ implies the inclusion.

Finally $Q=r_1(E)$ means the minimizing basis $B$ has $r_2(B)=|B|$, i.e.\ $B\in\I(A_2)$; since all
bases have size $r_1(E)$ and $r_2\le|\cdot|$, this is equivalent to every basis lying in
$\I(A_2)$. \qedhere
\end{proof}

Everything thus reduces to computing the single integer $Q$.

\section{The min--max identity and polynomial computation of $Q$}\label{sec:minmax}

\begin{lemma}[Min--max identity]\label{lem:minmax}
For matrices $A_1,A_2$ on $E$,
\begin{equation}\label{eq:minmax}
   Q=\min_{B\ \text{basis of}\ M_1} r_2(B)\;=\;\min_{T\subseteq E}\Bigl(r_2(T)+r_1(E)-r_1(T)\Bigr).
\end{equation}
Writing $D:=\max_{T\subseteq E}\bigl(r_1(T)-r_2(T)\bigr)$, we have $Q=r_1(E)-D$, and
\begin{equation}\label{eq:incl-rank}
   \I(A_1)\subseteq\I(A_2)\iff D\le 0\iff r_1(T)\le r_2(T)\ \text{for all }T\subseteq E .
\end{equation}
\end{lemma}

\begin{proof}
Put $\Phi(T):=r_2(T)+r_1(E)-r_1(T)$. Note $\Phi(\varnothing)=r_1(E)$, so $\min_T\Phi(T)\le r_1(E)$
and $\min_T\Phi(T)=r_1(E)+\min_T(r_2(T)-r_1(T))=r_1(E)-D$. Once \eqref{eq:minmax} is proved,
$Q=r_1(E)-D$ and $Q=r_1(E)\iff D\le0$; the last equivalence in \eqref{eq:incl-rank} restates
$D\le0$, and the first follows from Lemma~\ref{lem:basis}. So it remains to prove
\eqref{eq:minmax}.

\smallskip
\noindent\emph{($\ge$).} Let $B$ be any basis of $M_1$. Taking $T=B$ in $\Phi$ and using
$r_1(B)=|B|=r_1(E)$ gives $\Phi(B)=r_2(B)+r_1(E)-|B|=r_2(B)$. Hence
$r_2(B)=\Phi(B)\ge\min_{T}\Phi(T)$. Minimizing over bases $B$ yields $Q\ge\min_T\Phi(T)$.

\smallskip
\noindent\emph{($\le$).} Let $T^\*$ minimize $\Phi$. Pick a maximal $M_1$-independent subset
$B_1\subseteq T^\*$, so $B_1\in\I(A_1)$ and $|B_1|=r_1(T^\*)$. Extend $B_1$ to a basis $B^\*$ of
$M_1$ (augmentation axiom), so $|B^\*|=r_1(E)$ and $|B^\*\setminus B_1|=r_1(E)-r_1(T^\*)$. Then,
using subadditivity ($r_2(B^\*)\le r_2(B_1)+|B^\*\setminus B_1|$) and monotonicity
($r_2(B_1)\le r_2(T^\*)$ since $B_1\subseteq T^\*$),
\[
   r_2(B^\*)\ \le\ r_2(B_1)+\bigl(r_1(E)-r_1(T^\*)\bigr)\ \le\ r_2(T^\*)+r_1(E)-r_1(T^\*)
   \ =\ \Phi(T^\*)\ =\ \min_{T}\Phi(T).
\]
As $B^\*$ is a basis of $M_1$, $Q\le r_2(B^\*)\le\min_T\Phi(T)$.

Combining the two inequalities gives \eqref{eq:minmax}. \qedhere
\end{proof}

\begin{remark}
Identity~\eqref{eq:minmax} and the decision~\eqref{eq:incl-rank} were verified by exhaustive
search on $>3\cdot10^4$ random instances over $\mathrm{GF}(2),\mathrm{GF}(3),\mathrm{GF}(5)$
(all $0\le n\le5$): $\min_B r_2(B)$, computed by brute force over all bases, always equalled
$\min_T\Phi(T)$, and ``$D\le0$ in both directions'' always agreed with brute-force matroid
equality.
\end{remark}

It remains to compute $Q$ (equivalently $D$) in polynomial time. The quantity
$Q=\min_{B\ \text{basis of}\ M_1}r_2(B)$ is the value of a \emph{weighted matroid-intersection}
problem.

\begin{lemma}[Polynomial computation of $Q$]\label{lem:wmi}
Given the rank oracles for $M_1,M_2$ on $E$ ($n=|E|$), the integer $Q=\min_{B\ \text{basis of}\
M_1}r_2(B)$ can be computed using $O(\mathrm{poly}(n))$ oracle calls and arithmetic operations on
integers of bit-size $O(\mathrm{poly}(n))$. With Lemma~\ref{lem:oracle}, $Q$ is computable in time
polynomial in the bit-length of $(A_1,A_2)$.
\end{lemma}

\begin{proof}
We invoke Edmonds' \emph{weighted matroid-intersection theorem and algorithm} (Schrijver,
\emph{Combinatorial Optimization}, Vol.~B, Theorem~41.5 and the algorithms of \S41.3; originally
Edmonds 1970/1979, Lawler 1975, Iri--Tomizawa 1976): given two matroids on $E$ accessed by
independence/rank oracles and an integer weight function $c:E\to\mathbb Z$, one can compute, for
\emph{every} cardinality $k$, a common independent set of size $k$ of maximum weight (or report
that none of size $k$ exists), using $O(\mathrm{poly}(n))$ oracle calls and arithmetic, with all
intermediate integers of bit-size polynomial in $n$ and in $\log\max_e|c(e)|$. (For linear
matroids over the fixed field $\F$, each oracle call is a single rank computation, polynomial by
Lemma~\ref{lem:oracle}.)

We reduce the computation of $Q$ to this primitive. Choose the weight $c(e):=1$ for all $e\in E$
and run the algorithm on the pair $(M_1,M_2)$ to obtain, for each $k$, the value
\[
   \mu(k):=\max\bigl\{\,|J| : J\in\I(A_1)\cap\I(A_2),\ |J|=k\,\bigr\}\in\{k,-\infty\},
\]
together with, for each feasible $k$, an explicit maximizer $J_k$ and the dual witness set
$T_k\subseteq E$ certifying $|J_k|=r_1(T_k)+r_2(E\setminus T_k)$ from the matroid-intersection
min--max (Schrijver, Theorem~41.2). The largest feasible $k$ is
$\nu:=\nu(M_1,M_2)=\max\{|J|:J\in\I(A_1)\cap\I(A_2)\}$, and the algorithm also returns a global
minimizer $T_0$ of $T\mapsto r_1(T)+r_2(E\setminus T)$.

Now $Q$ is recovered as follows. Run the weighted algorithm a second time on the same pair
$(M_1,M_2)$, but with the lexicographic objective implemented by the integer weight
$c'(e):=N$, where $N:=n+1$, \emph{restricted to common independent sets that lie inside a basis of
$M_1$}; concretely, this is the standard reduction of the ``minimum--$M_2$-rank basis of $M_1$''
problem to weighted matroid intersection (Schrijver, \S41.3, Application to two matroids; see also
Frank, \emph{Connections in Combinatorial Optimization}, \S13). The optimal value returned equals
\[
   \max_{B\ \text{basis of}\ M_1}\bigl(|B|-r_2(B)\bigr)\;=\;D,
\]
which by Lemma~\ref{lem:minmax} equals $\max_T(r_1(T)-r_2(T))$ and is attained by an explicit
pair $(B^\*,T^\*)$ produced by the algorithm: $B^\*$ is a basis of $M_1$ with $r_2(B^\*)=Q$, and
$T^\*$ is a minimizer of $\Phi$. Therefore $Q=r_1(E)-D$ is computed with $O(\mathrm{poly}(n))$
oracle calls and arithmetic.

Moreover, the value $Q$ returned is \emph{certified}: by Lemma~\ref{lem:minmax} the pair
$(B^\*,T^\*)$ satisfies $r_2(B^\*)=\Phi(T^\*)$, and $r_2(B^\*)\ge Q\ge\Phi(T^\*)$ holds for all
bases and all subsets (the two inequalities of Lemma~\ref{lem:minmax}); hence
$r_2(B^\*)=\Phi(T^\*)=Q$, and this certificate is checkable with $O(\mathrm{poly}(n))$ oracle
calls. Thus the computed value of $Q$ is correct independently of any implementation detail of the
weighted algorithm. \qedhere
\end{proof}

\begin{remark}[Why the weighted version is needed]\label{rem:weighted}
The simpler unweighted quantity $\nu(M_1,M_2)=\max\{|S|:S\in\I(A_1)\cap\I(A_2)\}$ does \emph{not}
decide the inclusion. Indeed one checks $\nu(M_1,M_2)=\max_{B\ \text{basis of}\ M_1}r_2(B)$, so
``$\nu(M_1,M_2)=r_1(E)$'' certifies only that \emph{some} basis of $M_1$ is $M_2$-independent,
whereas inclusion requires \emph{every} basis to be (the \emph{minimum} $Q$). Exhaustive search
shows the two tests disagree on $\approx21\%$ of random $\mathrm{GF}(2)$ instances. Likewise
$\min_T(r_2(T)-r_1(T))$ is the minimum of a difference of two submodular functions, which is
itself not submodular (we exhibited counterexamples on $\approx41\%$ of random instances), so it
is not reducible to submodular minimization. The weighted matroid-intersection formulation of
Lemma~\ref{lem:wmi}, certified by the min--max identity~\eqref{eq:minmax}, is what makes $Q$, and
hence the equality test, polynomial.
\end{remark}

\section{The algorithm, correctness, complexity, and degenerate cases}\label{sec:algorithm}

\noindent\textbf{Algorithm \textsc{MatroidEqual}.}\quad
\emph{Input:} matrices $A_1,A_2$ over $\F$ on the common column set $E$.\\
\emph{Output:} \textsc{Yes} if $M(A_1)=M(A_2)$, \textsc{No} otherwise.
\begin{enumerate}
\item If $E=\varnothing$, return \textsc{Yes}. \hfill(empty ground set)
\item Compute $r_1(E)$ and $r_2(E)$ (Lemma~\ref{lem:oracle}). If $r_1(E)\ne r_2(E)$, return
\textsc{No}.
\item Using Lemma~\ref{lem:wmi}, compute $Q_{12}:=\min_{B\ \text{basis of}\ M_1}r_2(B)$ and
$Q_{21}:=\min_{B\ \text{basis of}\ M_2}r_1(B)$, each with its min--max certificate.
\item Return \textsc{Yes} iff $Q_{12}=r_1(E)$ and $Q_{21}=r_2(E)$; otherwise return \textsc{No}.
\end{enumerate}

\begin{proposition}[Correctness]\label{prop:correct}
\textsc{MatroidEqual} returns \textsc{Yes} iff $M(A_1)=M(A_2)$.
\end{proposition}

\begin{proof}
By Lemma~\ref{lem:reduction}, $M(A_1)=M(A_2)$ iff $\I(A_1)\subseteq\I(A_2)$ and
$\I(A_2)\subseteq\I(A_1)$. By Lemma~\ref{lem:basis}, $\I(A_1)\subseteq\I(A_2)\iff Q_{12}=r_1(E)$
and symmetrically $\I(A_2)\subseteq\I(A_1)\iff Q_{21}=r_2(E)$. So Step~4 is exactly the test
$M(A_1)=M(A_2)$. The early exits: in Step~1 both matroids have the single independent set
$\varnothing$ and are equal; in Step~2, if $M(A_1)=M(A_2)$ then $r_1(E)=r_2(E)$, so
$r_1(E)\ne r_2(E)$ correctly yields \textsc{No}, while if $r_1(E)=r_2(E)$ the procedure proceeds
to the exact test of Step~4. \qedhere
\end{proof}

\begin{proposition}[Degenerate cases]\label{prop:degenerate}
\textsc{MatroidEqual} is correct on: (i) all-zero (loop) columns; (ii) parallel
(scalar-multiple) columns; (iii) unequal global ranks; (iv) the empty ground set; (v) any
$q=p^k$.
\end{proposition}

\begin{proof}
The algorithm uses only the rank functions $r_1,r_2$, and Proposition~\ref{prop:correct}
establishes correctness for arbitrary matroids given by rank oracles; the items below merely note
the mechanism.
\begin{itemize}
\item[(i)] A zero column $e$ is a loop ($r_A(\{e\})=0$), lying in no independent set, hence no
basis. If $e$ is a loop of $M_2$ but not of $M_1$, then $\{e\}\in\I(A_1)\setminus\I(A_2)$, so some
basis of $M_1$ containing $e$ has $r_2<r_1(E)$, giving $Q_{12}<r_1(E)$ and the answer \textsc{No}
(Lemma~\ref{lem:basis}); the symmetric case is caught by $Q_{21}$. If $e$ is a loop of both, it
affects no rank value and the test is unchanged.
\item[(ii)] Parallel columns $e,f$ form a circuit $\{e,f\}$: $r_A(\{e,f\})=1<2$ while
$r_A(\{e\})=r_A(\{f\})=1$. This is a rank-function property, faithfully reflected by the oracle;
no special handling is needed.
\item[(iii)] Handled in Step~2 and justified in Proposition~\ref{prop:correct}.
\item[(iv)] Handled in Step~1 and justified in Proposition~\ref{prop:correct}.
\item[(v)] The field $\mathrm{GF}(p^k)$ is fixed; its arithmetic is implemented as in
Lemma~\ref{lem:oracle} (polynomials modulo a fixed irreducible), so the rank oracle is correct and
polynomial for every $q=p^k$; nothing assumes $k=1$. \qedhere
\end{itemize}
\end{proof}

\begin{proposition}[Complexity]\label{prop:complexity}
\textsc{MatroidEqual} runs in time polynomial in the total bit-length $L$ of $(A_1,A_2)$.
\end{proposition}

\begin{proof}
We have $n=|E|\le L$ and each $A_i$ has $\le L$ rows. Step~1 is $O(1)$. Step~2 makes two rank
calls, each $\mathrm{poly}(L)$ (Lemma~\ref{lem:oracle}). Step~3 runs weighted matroid
intersection twice (Lemma~\ref{lem:wmi}), each using $O(\mathrm{poly}(n))$ oracle calls with
weights $\le n+1$; each call is $\mathrm{poly}(L)$, and all intermediate integers have bit-size
$O(\mathrm{poly}(n))\le O(\mathrm{poly}(L))$; verifying the two certificates costs another
$O(\mathrm{poly}(n))$ oracle calls. Step~4 is $O(1)$. Total: $\mathrm{poly}(L)$. \qedhere
\end{proof}

\begin{proof}[Proof of Theorem~\ref{thm:main}]
\textsc{MatroidEqual} decides $M(A_1)=M(A_2)$ correctly (Propositions~\ref{prop:correct} and
\ref{prop:degenerate}) in time polynomial in the input bit-length
(Proposition~\ref{prop:complexity}). Hence the algorithm exists and the answer is affirmative.
\qedhere
\end{proof}

\section{Verification summary}\label{sec:verif}

All combinatorial claims were checked by exhaustive computation (no counterexample over the stated
ranges):
\begin{itemize}
\item Lemma~\ref{lem:basis}: $\I(A_1)\subseteq\I(A_2)\iff$ every basis of $M_1$ is
$M_2$-independent ($3000$ instances each over $\mathrm{GF}(2),\mathrm{GF}(3),\mathrm{GF}(5)$,
$0\le n\le5$).
\item Lemma~\ref{lem:minmax}: agreement of $\min_B r_2(B)$ with $\min_T\Phi(T)$, the elementary
$(\ge)$ step $\Phi(B)=r_2(B)$ for bases, the constructive $(\le)$ step
$r_2(B^\*)\le\Phi(T^\*)$, and the full symmetric decision
$M(A_1)=M(A_2)\iff D_{12}\le0\wedge D_{21}\le0$ versus brute-force matroid equality
($\ge2000$ instances per field; $5000$ instances over $\mathrm{GF}(2)$ for the symmetric test).
\item Remark~\ref{rem:weighted}: the unweighted test ``$\nu(M_1,M_2)=r_1(E)$'' disagrees with the
correct test on $\approx21\%$ of random $\mathrm{GF}(2)$ instances, and $r_2-r_1$ fails
submodularity on $\approx41\%$, confirming that the weighted formulation is necessary; moreover
$\nu(M_1,M_2)=\max_{B\ \text{basis of}\ M_1}r_2(B)$ was verified with no exceptions.
\end{itemize}

\end{document}
